%=====================================
%   geometry.tex
%   幾何学まとめノート toshi2019
%=====================================

% -----------------------
% preamble
% -----------------------
% ここから本文 (\begin{document}) までの
% ソースコードに変更を加えた場合は
% 編集者まで連絡してください. 
% Don't change preamble code yourself. 
% If you add something
% (usepackage, newtheorem, newcommand, renewcommand),
% please tell it 
% to the editor of institutional paper of RUMS.

% ------------------------
% documentclass
% ------------------------
\documentclass[9pt, a4paper, dvipdfmx]{jsarticle}

% ------------------------
% usepackage
% ------------------------
\usepackage{algorithm}
\usepackage{algorithmic}
\usepackage{amscd}
\usepackage{amsfonts}
\usepackage{amsmath}
\usepackage[psamsfonts]{amssymb}
\usepackage{amsthm}
\usepackage{ascmac}
\usepackage{bm}
\usepackage{color}
\usepackage{enumerate}
\usepackage{fancybox}
\usepackage[stable]{footmisc}
\usepackage{graphicx}
\usepackage{listings}
\usepackage{mathrsfs}
\usepackage{mathtools}
\usepackage{otf}
\usepackage{pifont}
\usepackage{proof}
\usepackage{subfigure}
\usepackage{tikz}
\usepackage{verbatim}
\usepackage[all]{xy}

\usetikzlibrary{cd}
\usetikzlibrary{arrows.meta}


% ================================
% パッケージを追加する場合のスペース 
\usepackage[dvipdfmx]{hyperref}
\usepackage{xcolor}
\definecolor{darkgreen}{rgb}{0,0.45,0} 
\definecolor{darkred}{rgb}{0.75,0,0}
\definecolor{darkblue}{rgb}{0,0,0.6} 
\hypersetup{
    colorlinks=true,
    citecolor=darkgreen,
    linkcolor=darkred,
    urlcolor=darkblue,
}
\usepackage{pxjahyper}

%\usepackage[mathscr]{euscript} % mathscr のフォントを変更

%\usepackage{mmasym}

%=================================


% --------------------------
% theoremstyle
% --------------------------
\theoremstyle{definition}

% --------------------------
% newtheoem
% --------------------------

% 日本語で定理, 命題, 証明などを番号付きで用いるためのコマンドです. 
% If you want to use theorem environment in Japanece, 
% you can use these code. 
% Attention!
% All theorem enivironment numbers depend on 
% only section numbers.
\newtheorem{Axiom}{公理}[section]
\newtheorem{Definition}[Axiom]{定義}
\newtheorem{Theorem}[Axiom]{定理}
\newtheorem{Proposition}[Axiom]{命題}
\newtheorem{Lemma}[Axiom]{補題}
\newtheorem{Corollary}[Axiom]{系}
\newtheorem{Example}[Axiom]{例}
\newtheorem{Claim}[Axiom]{主張}
\newtheorem{Property}[Axiom]{性質}
\newtheorem{Attention}[Axiom]{注意}
\newtheorem{Question}[Axiom]{問}
\newtheorem{Problem}[Axiom]{問題}
\newtheorem{Consideration}[Axiom]{考察}
\newtheorem{Alert}[Axiom]{警告}
\newtheorem{Fact}[Axiom]{事実}


% 日本語で定理, 命題, 証明などを番号なしで用いるためのコマンドです. 
% If you want to use theorem environment with no number in Japanese, You can use these code.
\newtheorem*{Axiom*}{公理}
\newtheorem*{Definition*}{定義}
\newtheorem*{Theorem*}{定理}
\newtheorem*{Proposition*}{命題}
\newtheorem*{Lemma*}{補題}
\newtheorem*{Example*}{例}
\newtheorem*{Corollary*}{系}
\newtheorem*{Claim*}{主張}
\newtheorem*{Property*}{性質}
\newtheorem*{Attention*}{注意}
\newtheorem*{Question*}{問}
\newtheorem*{Problem*}{問題}
\newtheorem*{Consideration*}{考察}
\newtheorem*{Alert*}{警告}
\newtheorem{Fact*}{事実}


% 英語で定理, 命題, 証明などを番号付きで用いるためのコマンドです. 
%　If you want to use theorem environment in English, You can use these code.
%all theorem enivironment number depend on only section number.
\newtheorem{Axiom+}{Axiom}[section]
\newtheorem{Definition+}[Axiom+]{Definition}
\newtheorem{Theorem+}[Axiom+]{Theorem}
\newtheorem{Proposition+}[Axiom+]{Proposition}
\newtheorem{Lemma+}[Axiom+]{Lemma}
\newtheorem{Example+}[Axiom+]{Example}
\newtheorem{Corollary+}[Axiom+]{Corollary}
\newtheorem{Claim+}[Axiom+]{Claim}
\newtheorem{Property+}[Axiom+]{Property}
\newtheorem{Attention+}[Axiom+]{Attention}
\newtheorem{Question+}[Axiom+]{Question}
\newtheorem{Problem+}[Axiom+]{Problem}
\newtheorem{Consideration+}[Axiom+]{Consideration}
\newtheorem{Alert+}{Alert}
\newtheorem{Fact+}[Axiom+]{Fact}
\newtheorem{Remark+}[Axiom+]{Remark}

% ----------------------------
% commmand
% ----------------------------
% 執筆に便利なコマンド集です. 
% コマンドを追加する場合は下のスペースへ. 

% 集合の記号 (黒板文字)
\newcommand{\NN}{\mathbb{N}}
\newcommand{\ZZ}{\mathbb{Z}}
\newcommand{\QQ}{\mathbb{Q}}
\newcommand{\RR}{\mathbb{R}}
\newcommand{\CC}{\mathbb{C}}
\newcommand{\PP}{\mathbb{P}}
\newcommand{\KK}{\mathbb{K}}


% 集合の記号 (太文字)
\newcommand{\nn}{\mathbf{N}}
\newcommand{\zz}{\mathbf{Z}}
\newcommand{\qq}{\mathbf{Q}}
\newcommand{\rr}{\mathbf{R}}
\newcommand{\cc}{\mathbf{C}}
\newcommand{\pp}{\mathbf{P}}
\newcommand{\kk}{\mathbf{k}}

% 特殊な写像の記号
\newcommand{\ev}{\mathop{\mathrm{ev}}\nolimits} % 値写像
\newcommand{\pr}{\mathop{\mathrm{pr}}\nolimits} % 射影

% スクリプト体にするコマンド
%   例えば　{\mcal C} のように用いる
\newcommand{\mcal}{\mathcal}

% 花文字にするコマンド 
%   例えば　{\h C} のように用いる
\newcommand{\h}{\mathscr}

% ヒルベルト空間などの記号
\newcommand{\F}{\mcal{F}}
\newcommand{\X}{\mcal{X}}
\newcommand{\Y}{\mcal{Y}}
\newcommand{\Hil}{\mcal{H}}
\newcommand{\RKHS}{\Hil_{k}}
\newcommand{\Loss}{\mcal{L}_{D}}
\newcommand{\MLsp}{(\X, \Y, D, \Hil, \Loss)}

% 偏微分作用素の記号
\newcommand{\p}{\partial}

% 角カッコの記号 (内積は下にマクロがあります)
\newcommand{\lan}{\langle}
\newcommand{\ran}{\rangle}



% 圏の記号など
\newcommand{\Set}{{\bf Set}}
\newcommand{\Vect}{{\bf Vect}}
\newcommand{\FDVect}{{\bf FDVect}}
%\newcommand{\Ring}{{\bf Ring}}
\newcommand{\Ab}{{\bf Ab}}
\newcommand{\Mod}{\mathop{\mathrm{Mod}}\nolimits}
\newcommand{\Modf}{\mathop{\mathrm{Mod}^\mathrm{f}}\nolimits}
\newcommand{\CGA}{{\bf CGA}}
\newcommand{\GVect}{{\bf GVect}}
\newcommand{\Lie}{{\bf Lie}}
\newcommand{\dLie}{{\bf Liec}}



% 射の集合など
\newcommand{\Map}{\mathop{\mathrm{Map}}\nolimits} % 写像の集合
\newcommand{\Hom}{\mathop{\mathrm{Hom}}\nolimits} % 射集合
\newcommand{\End}{\mathop{\mathrm{End}}\nolimits} % 自己準同型の集合
\newcommand{\Aut}{\mathop{\mathrm{Aut}}\nolimits} % 自己同型の集合
\newcommand{\Mor}{\mathop{\mathrm{Mor}}\nolimits} % 射集合
\newcommand{\Ker}{\mathop{\mathrm{Ker}}\nolimits} % 核
\newcommand{\Img}{\mathop{\mathrm{Im}}\nolimits} % 像
\newcommand{\Cok}{\mathop{\mathrm{Coker}}\nolimits} % 余核
\newcommand{\Cim}{\mathop{\mathrm{Coim}}\nolimits} % 余像

% その他便利なコマンド
\newcommand{\dip}{\displaystyle} % 本文中で数式モード
\newcommand{\e}{\varepsilon} % イプシロン
\newcommand{\dl}{\delta} % デルタ
\newcommand{\pphi}{\varphi} % ファイ
\newcommand{\ti}{\tilde} % チルダ
\newcommand{\pal}{\parallel} % 平行
\newcommand{\op}{{\rm op}} % 双対を取る記号
\newcommand{\lcm}{\mathop{\mathrm{lcm}}\nolimits} % 最小公倍数の記号
\newcommand{\Probsp}{(\Omega, \F, \P)} 
\newcommand{\argmax}{\mathop{\rm arg~max}\limits}
\newcommand{\argmin}{\mathop{\rm arg~min}\limits}





% ================================
% コマンドを追加する場合のスペース 
%\newcommand{\OO}{\mcal{O}}



\makeatletter
\renewenvironment{proof}[1][\proofname]{\par
  \pushQED{\qed}%
  \normalfont \topsep6\p@\@plus6\p@\relax
  \trivlist
  \item[\hskip\labelsep
%        \itshape
         \bfseries
%    #1\@addpunct{.}]\ignorespaces
    {#1}]\ignorespaces
}{%
  \popQED\endtrivlist\@endpefalse
}
\makeatother

\renewcommand{\proofname}{証明.}



%\renewcommand\proofname{\bf 証明} % 証明
\numberwithin{equation}{section}
\newcommand{\cTop}{\textsf{Top}}
%\newcommand{\cOpen}{\textsf{Open}}
\newcommand{\Op}{\mathop{\textsf{Op}}\nolimits}
\newcommand{\Ob}{\mathop{\textrm{Ob}}\nolimits}
\newcommand{\id}{\mathop{\mathrm{id}}\nolimits}
\newcommand{\pt}{\mathop{\mathrm{pt}}\nolimits}
\newcommand{\res}{\mathop{\rho}\nolimits}
\newcommand{\A}{\mcal{A}}
\newcommand{\B}{\mcal{B}}
\newcommand{\C}{\mcal{C}}
\newcommand{\D}{\mcal{D}}
\newcommand{\E}{\mcal{E}}
\newcommand{\G}{\mcal{G}}
%\newcommand{\H}{\mcal{H}}
\newcommand{\I}{\mcal{I}}
\newcommand{\J}{\mcal{J}}
\newcommand{\OO}{\mcal{O}}
\newcommand{\Ring}{\mathop{\textsf{Ring}}\nolimits}
\newcommand{\cAb}{\mathop{\textsf{Ab}}\nolimits}
%\newcommand{\Ker}{\mathop{\mathrm{Ker}}\nolimits}
\newcommand{\im}{\mathop{\mathrm{Im}}\nolimits}
\newcommand{\Coker}{\mathop{\mathrm{Coker}}\nolimits}
\newcommand{\Coim}{\mathop{\mathrm{Coim}}\nolimits}
\newcommand{\rank}{\mathop{\mathrm{rank}}\nolimits}
\newcommand{\Ht}{\mathop{\mathrm{Ht}}\nolimits}
\newcommand{\supp}{\mathop{\mathrm{supp}}\nolimits}
\newcommand{\colim}{\mathop{\mathrm{colim}}}
\newcommand{\Tor}{\mathop{\mathrm{Tor}}\nolimits}

\newcommand{\cat}{\mathscr{C}}

%筆記体
\newcommand{\cA}{\mcal{A}}
\newcommand{\cB}{\mcal{B}}
\newcommand{\cC}{\mcal{C}}
\newcommand{\cD}{\mcal{D}}
\newcommand{\cE}{\mcal{E}}
\newcommand{\cF}{\mcal{F}}
\newcommand{\cG}{\mcal{G}}
\newcommand{\cH}{\mcal{H}}
\newcommand{\cI}{\mcal{I}}
\newcommand{\cJ}{\mcal{J}}
\newcommand{\cK}{\mcal{K}}
\newcommand{\cL}{\mcal{L}}
\newcommand{\cM}{\mcal{M}}
\newcommand{\cN}{\mcal{N}}
\newcommand{\cO}{\mcal{O}}
\newcommand{\cP}{\mcal{P}}
\newcommand{\cQ}{\mcal{Q}}
\newcommand{\cR}{\mcal{R}}
\newcommand{\cS}{\mcal{S}}
\newcommand{\cT}{\mcal{T}}
\newcommand{\cU}{\mcal{U}}
\newcommand{\cV}{\mcal{V}}
\newcommand{\cW}{\mcal{W}}
\newcommand{\cX}{\mcal{X}}
\newcommand{\cY}{\mcal{Y}}
\newcommand{\cZ}{\mcal{Z}}


\newcommand{\scA}{\mathscr{A}}
\newcommand{\scB}{\mathscr{B}}
\newcommand{\scC}{\mathscr{C}}
\newcommand{\scD}{\mathscr{D}}
\newcommand{\scE}{\mathscr{E}}
\newcommand{\scF}{\mathscr{F}}
\newcommand{\scN}{\mathscr{N}}
\newcommand{\scO}{\mathscr{O}}
\newcommand{\scV}{\mathscr{V}}
\newcommand{\scU}{\mathscr{U}}


\newcommand{\ibA}{\mathop{\text{\textit{\textbf{A}}}}}
\newcommand{\ibB}{\mathop{\text{\textit{\textbf{B}}}}}
\newcommand{\ibC}{\mathop{\text{\textit{\textbf{C}}}}}
\newcommand{\ibD}{\mathop{\text{\textit{\textbf{D}}}}}
\newcommand{\ibE}{\mathop{\text{\textit{\textbf{E}}}}}
\newcommand{\ibF}{\mathop{\text{\textit{\textbf{F}}}}}
\newcommand{\ibG}{\mathop{\text{\textit{\textbf{G}}}}}
\newcommand{\ibH}{\mathop{\text{\textit{\textbf{H}}}}}
\newcommand{\ibI}{\mathop{\text{\textit{\textbf{I}}}}}
\newcommand{\ibJ}{\mathop{\text{\textit{\textbf{J}}}}}
\newcommand{\ibK}{\mathop{\text{\textit{\textbf{K}}}}}
\newcommand{\ibL}{\mathop{\text{\textit{\textbf{L}}}}}
\newcommand{\ibM}{\mathop{\text{\textit{\textbf{M}}}}}
\newcommand{\ibN}{\mathop{\text{\textit{\textbf{N}}}}}
\newcommand{\ibO}{\mathop{\text{\textit{\textbf{O}}}}}
\newcommand{\ibP}{\mathop{\text{\textit{\textbf{P}}}}}
\newcommand{\ibQ}{\mathop{\text{\textit{\textbf{Q}}}}}
\newcommand{\ibR}{\mathop{\text{\textit{\textbf{R}}}}}
\newcommand{\ibS}{\mathop{\text{\textit{\textbf{S}}}}}
\newcommand{\ibT}{\mathop{\text{\textit{\textbf{T}}}}}
\newcommand{\ibU}{\mathop{\text{\textit{\textbf{U}}}}}
\newcommand{\ibV}{\mathop{\text{\textit{\textbf{V}}}}}
\newcommand{\ibW}{\mathop{\text{\textit{\textbf{W}}}}}
\newcommand{\ibX}{\mathop{\text{\textit{\textbf{X}}}}}
\newcommand{\ibY}{\mathop{\text{\textit{\textbf{Y}}}}}
\newcommand{\ibZ}{\mathop{\text{\textit{\textbf{Z}}}}}

\newcommand{\ibx}{\mathop{\text{\textit{\textbf{x}}}}}

%\newcommand{\Comp}{\mathop{\mathrm{C}}\nolimits}
%\newcommand{\Komp}{\mathop{\mathrm{K}}\nolimits}
%\newcommand{\Domp}{\mathop{\mathsf{D}}\nolimits}%複体のホモトピー圏
%\newcommand{\Comp}{\mathrm{C}}
%\newcommand{\Komp}{\mathrm{K}}
%\newcommand{\Domp}{\mathsf{D}}%複体のホモトピー圏

\newcommand{\Comp}{\mathop{\mathrm{C}}\nolimits}
\newcommand{\Komp}{\mathop{\mathsf{K}}\nolimits}
\newcommand{\Domp}{\mathord{\mathsf{D}}}%\nolimits}
\newcommand{\Kompl}{\mathop{\mathsf{K}^\mathrm{+}}\nolimits}
\newcommand{\Kompu}{\mathop{\mathsf{K}^\mathrm{-}}\nolimits}
\newcommand{\Kompb}{\mathop{\mathsf{K}^\mathrm{b}}\nolimits}
\newcommand{\Dompl}{\mathop{\mathsf{D}^\mathrm{+}}\nolimits}
\newcommand{\Dompu}{\mathop{\mathsf{D}^\mathrm{-}}\nolimits}
\newcommand{\Dompb}{\mathop{\mathsf{D}^\mathrm{b}}\nolimits}
\newcommand{\Dompbf}{\mathop{\mathsf{D}_\mathrm{f}^\mathrm{b}}\nolimits}
\newcommand{\Domplb}{\mathop{\mathsf{D}^\mathrm{lb}}\nolimits}




\newcommand{\CCat}{\Comp(\cat)}
\newcommand{\KCat}{\Komp(\cat)}
\newcommand{\DCat}{\Domp(\cat)}%圏Cの複体のホモトピー圏
\newcommand{\HOM}{\mathop{\mathscr{H}\hspace{-2pt}om}\nolimits}%内部Hom
\newcommand{\RHOM}{\mathop{\mathrm{R}\hspace{-1.5pt}\HOM}\nolimits}

\newcommand{\muS}{\mathop{\mathrm{SS}}\nolimits}
\newcommand{\RG}{\mathop{\mathrm{R}\hspace{-0pt}\Gamma}\nolimits}
\newcommand{\RHom}{\mathop{\mathrm{R}\hspace{-1.5pt}\Hom}\nolimits}
\newcommand{\Rder}{\mathrm{R}}

\newcommand{\simar}{\mathrel{\overset{\sim}{\rightarrow}}}%同型右矢印
\newcommand{\simarr}{\mathrel{\overset{\sim}{\longrightarrow}}}%同型右矢印
\newcommand{\simra}{\mathrel{\overset{\sim}{\leftarrow}}}%同型左矢印
\newcommand{\simrra}{\mathrel{\overset{\sim}{\longleftarrow}}}%同型左矢印

\newcommand{\hocolim}{{\mathrm{hocolim}}}
\newcommand{\indlim}[1][]{\mathop{\varinjlim}\limits_{#1}}
\newcommand{\sindlim}[1][]{\smash{\mathop{\varinjlim}\limits_{#1}}\,}
\newcommand{\Pro}{\mathrm{Pro}}
\newcommand{\Ind}{\mathrm{Ind}}
\newcommand{\prolim}[1][]{\mathop{\varprojlim}\limits_{#1}}
\newcommand{\sprolim}[1][]{\smash{\mathop{\varprojlim}\limits_{#1}}\,}

\newcommand{\Sh}{\mathrm{Sh}}
\newcommand{\PSh}{\mathrm{PSh}}

\newcommand{\rmD}{\mathrm{D}}

\newcommand{\Lloc}[1][]{\mathord{\mathcal{L}^1_{\mathrm{loc},{#1}}}}
\newcommand{\ori}{\mathord{\mathrm{or}}}
\newcommand{\Db}{\mathord{\mathscr{D}b}}

\newcommand{\codim}{\mathop{\mathrm{codim}}\nolimits}



\newcommand{\gld}{\mathop{\mathrm{gld}}\nolimits}
\newcommand{\wgld}{\mathop{\mathrm{wgld}}\nolimits}


\newcommand{\tens}[1][]{\mathbin{\otimes_{\raise1.5ex\hbox to-.1em{}{#1}}}}
\newcommand{\etens}{\mathbin{\boxtimes}}
\newcommand{\ltens}[1][]{\mathbin{\overset{\mathrm{L}}\tens}_{#1}}
\newcommand{\mtens}[1][]{\mathbin{\overset{\mathrm{\mu}}\tens}_{#1}}
\newcommand{\lltens}[1][]{{\mathop{\tens}\limits^{\mathrm{L}}_{#1}}}
\newcommand{\letens}{\overset{\mathrm{L}}{\etens}}
\newcommand{\detens}{\underline{\etens}}
\newcommand{\ldetens}{\overset{\mathrm{L}}{\underline{\etens}}}
\newcommand{\dtens}[1][]{{\overset{\mathrm{L}}{\underline{\otimes}}}_{#1}}

\newcommand{\blk}{\mathord{\ \cdot\ }}
\newcommand{\mres}[2][]{{\left.{#1}\right\rvert}_{#2}}


%\newcommand{\hocolim}{{\mathrm{hocolim}}}
%\newcommand{\indlim}[1][]{\mathop{\varinjlim}\limits_{#1}}
%\newcommand{\sindlim}[1][]{\smash{\mathop{\varinjlim}\limits_{#1}}\,}
%\newcommand{\Pro}{\mathrm{Pro}}
%\newcommand{\Ind}{\mathrm{Ind}}
%\newcommand{\prolim}[1][]{\mathop{\varprojlim}\limits_{#1}}
%\newcommand{\sprolim}[1][]{\smash{\mathop{\varprojlim}\limits_{#1}}\,}
\newcommand{\proolim}[1][]{\mathop{\text{\rm``{$\varprojlim$}''}}\limits_{#1}}
\newcommand{\sproolim}[1][]{\smash{\mathop{\rm``{\varprojlim}''}\limits_{#1}}}
\newcommand{\inddlim}[1][]{\mathop{\text{\rm``{$\varinjlim$}''}}\limits_{#1}}
\newcommand{\sinddlim}[1][]{\smash{\mathop{\text{\rm``{$\varinjlim$}''}}\limits_{#1}}\,}
\newcommand{\ooplus}{\mathop{\text{\rm``{$\oplus$}''}}\limits}
\newcommand{\bbigsqcup}{\mathop{``\bigsqcup''}\limits}
\newcommand{\bsqcup}{\mathop{``\sqcup''}\limits}
\newcommand{\dsum}[1][]{\mathbin{\oplus_{#1}}}

\newcommand{\Fct}{\mathop{\mathsf{Fct}}\nolimits}

\newcommand{\pb}[1][]{\mathbin{\mathop{\times}\limits_{#1}}}





%================================================
% 自前の定理環境
%   https://mathlandscape.com/latex-amsthm/
% を参考にした
\newtheoremstyle{mystyle}%   % スタイル名
    {5pt}%                   % 上部スペース
    {5pt}%                   % 下部スペース
    {}%              % 本文フォント
    {}%                  % 1行目のインデント量
    {\bfseries}%                      % 見出しフォント
    {.}%                     % 見出し後の句読点
    {12pt}%                     % 見出し後のスペース
    {\thmname{#1}\thmnumber{ #2}\thmnote{{\hspace{2pt}\normalfont (#3)}}}% % 見出しの書式

\theoremstyle{mystyle}
\newtheorem{AXM}{公理}[section]
\newtheorem{DFN}[AXM]{定義}
\newtheorem{THM}[AXM]{定理}
\newtheorem*{THM*}{定理}
\newtheorem{PRP}[AXM]{命題}
\newtheorem{LMM}[AXM]{補題}
\newtheorem{CRL}[AXM]{系}
\newtheorem{EG}[AXM]{例}
\newtheorem*{EG*}{例}
\newtheorem{RMK}[AXM]{注意}
\newtheorem{CNV}[AXM]{約束}
\newtheorem{CMT}[AXM]{コメント}
\newtheorem*{CMT*}{コメント}
\newtheorem{NTN}[AXM]{記号}

% 定理環境ここまで
%====================================================

\usepackage{framed}
\definecolor{lightgray}{rgb}{0.75,0.75,0.75}
\renewenvironment{leftbar}{%
  \def\FrameCommand{\textcolor{lightgray}{\vrule width 4pt} \hspace{10pt}}% 
  \MakeFramed {\advance\hsize-\width \FrameRestore}}%
{\endMakeFramed}
\newenvironment{redleftbar}{%
  \def\FrameCommand{\textcolor{lightgray}{\vrule width 1pt} \hspace{10pt}}% 
  \MakeFramed {\advance\hsize-\width \FrameRestore}}%
 {\endMakeFramed}



\renewcommand{\Re}{\mathop{\mathrm{Re}}\nolimits}

% =================================





% ---------------------------
% new definition macro
% ---------------------------
% 便利なマクロ集です

% 内積のマクロ
%   例えば \inner<\pphi | \psi> のように用いる
\def\inner<#1>{\langle #1 \rangle}

% ================================
% マクロを追加する場合のスペース 

%=================================



% 位相空間
\newcommand{\Int}[1][]{\mathop{\mathrm{Int}}\nolimits_{#1}}
\newcommand{\Cl}[1][]{\mathop{\mathrm{Cl}}\nolimits_{#1}}
\newcommand{\Bd}[1][]{\mathop{\mathrm{Bd}}\nolimits_{#1}}
\newcommand{\bd}{\mathord{\partial}}
\newcommand{\Fr}[1][]{\mathop{\mathrm{Fr}}\nolimits_{#1}}


\newcommand{\mtan}[1][]{T\hspace{-1.75pt}{#1}}
\newcommand{\mtp}[1][]{{}^{t\!}{#1}} % transpose of the morphism
%\newcommand{\mpb}[1][]{{}^{t\!}{#1}'} % pull-back of the morphism
\newcommand{\mpb}[1][]{{#1}_{\pi}} % pull-back of the morphism
\newcommand{\mdiff}[1][]{{#1}_{d}} % pull-back of the morphism
\newcommand{\mptan}[1][]{{#1}_{\tau}} % pull-back of the morphism


%\newcommand{\pb}[1][]{\mathbin{\mathop{\times}\limits_{#1}}}

\newcommand{\cotb}[1][]{T^{\ast}{#1}}
\newcommand{\conm}[2][]{T_{#1}^{\hspace{0.7pt}\ast}{#2}}
\newcommand{\norb}[2][]{T_{#1}{#2}}

\newcommand{\cotz}[1][]{T^{\ast}_{#1}{#1}}

\newcommand{\cotn}[1][]{\dot{T}^{\ast}{#1}}


\newcommand{\muhom}{\mu hom}
\newcommand{\muHom}[1][]{\mathrm{Hom}^\mu_{\raise1.5ex\hbox to.1em{}#1}}

\newcommand{\mucat}[2][]{\mathsf{D}^{\mathrm{b}}_{#1}(#2)}
%\newcommand{\mucat}[2][]{\mathsf{D}^{\mathrm{b}}_{#1}(#2)}

\newcommand{\conv}[1][]{\mathbin{\mathop{\circ}\limits_{#1}}}

% ----------------------------
% documenet 
% ----------------------------
% 以下, 本文の執筆スペースです. 
% Your main code must be written between 
% begin document and end document.
% ---------------------------

\title{ハミルトンイソトピーの層量子化と分離不能性問題への応用}
\author{Toshi2019}
\date{\today\footnote{2025/02/24作成開始}}
\begin{document}
\maketitle
%\tableofcontents
\tableofcontents
\section*{はじめに}
\cite{GKS12}のまとめノート．ほぼ和訳？
\subsection*{注意}
ArXiv版とDuke版で節の番号がズレている．
このノートではDuke版に合わせる．
\section*{要約}
\begin{itemize}
    \item \(I\)：\(0\)を含む開区間
    \item \(M\)：実多様体
    \item \(\cotn[M]=\cotb[M]\smallsetminus\cotz[M]\)
    \item \(\varPhi=\left(\varphi_{t}\right)_{t\in I}\)：\(\varphi_0=\id_{\cotn[M]}\)をみたす同次シンプレクティックイソトピー
    \item \(\Lambda\subset\cotn[M]\times\cotn[M]\times\cotn[I]\)：\(\varPhi\)から定まる錐形ラグランジュ部分多様体
\end{itemize}
のとき，次の主張を示す．
\begin{quote}
    \(M\times M\times I\)上の層\(K\)で，
    その超局所台が\(\Lambda\)と零切断の合併に含まれ，
    \(t=0\)への制限が\(M\times M\)上の定数層となる
    ものがただ一つ存在する．
\end{quote}
分離不能性問題への応用として
\begin{quote}
    強いモース不等式がハミルトンイソトピーで安定であること
    
    接触幾何の設定での非負イソトピーに対する分離不能性に関する結果
\end{quote}
が得られる．

\section*{序論}

\begin{itemize}
    \item 1982--90：柏原--Schapira，層の超局所理論を導入，展開．中心的なアイデアは層の超局所台．
    \item \(M\)：\(C^\infty\)多様体
    \item \(\kk\)：1を持つ大域次元が有限な可換環
    \item \(\Dompb(\kk_M)\)：\(M\)上の\(\kk\)加群の層の有界導来圏
\end{itemize}
とする．\(F\in\Dompb(\kk_M)\)に対し，\(F\)の超局所台
\[
    \muS(F)\subset\cotb[M]
\]
が定まる．これは
\begin{itemize}
    \item 錐形閉集合，つまりファイバーごとに\(\rr^{+}\)作用で閉じた閉集合
    \item 余等方的
\end{itemize}
である．したがって，この理論は「錐形」．つまり
\begin{itemize}
    \item \(\rr^{+}\)作用に関して不変
    \item \(\cotb[M]\)のシンプレクティック構造でなく，同次シンプレクティック構造の方に関連している
\end{itemize}
そこで，非同次シンプレクティックの方を考えるには，変数を追加する．
この手法は，複素シンプレクティックの方ではPolleselo--Schapiraが，
実シンプレクティックの方ではTamarkinがすでに用いている．

非同次の理論を同次に帰着してから超局所理論を用いる．

逆は常にはできない．

\section{柏原--Schapira{\cite{KS90}}による層の超局所理論}
ここでは実多様体\(M\)として\(C^\infty\)級のものを考える．
\subsection{幾何的諸概念}
本論文において，\(C^1\)級写像\(
    f\colon M\to N
\)が\textbf{平滑} (smooth) であるとは，
その微分\(
    d_{x}f\colon T_{x}M\to T_{f(x)}N
\)が任意の\(x\in M\)に対して全射であることをいう．


\begin{equation}
    \vcenter{\xymatrix@C=36pt@R=36pt{
    TM
    \ar[d]^-{\tau_M}
    \ar[r]^-{f'}
    &
    M\pb[N] TN
    %\ar@{}[dr]|\square
    \ar[d]_-{}
    \ar[r]^-{\mptan[f]}
    &
    \cotb[N]
    \ar[d]^-{\tau_N}
    \\
    M
    \ar@{=}[r]
    &
    M
    \ar[r]^-{f}
    %\ar@{}[dr]|\square
    &
    N.
  }}
\end{equation}
双対的に，次の図式を得る．
\begin{equation}
    \vcenter{\xymatrix@C=36pt@R=36pt{
    \cotb[M]
    \ar[d]^-{\pi_M}
    &
    M\pb[N] \cotb[N]
    %\ar@{}[dr]|\square
    \ar[d]_-{}
    \ar[r]^-{\mpb[f]}
    \ar[l]_-{\mdiff[f]}
    &
    \cotb[N]
    \ar[d]^-{\pi_N}
    \\
    M
    \ar@{=}[r]
    &
    M
    \ar[r]^-{f}
    %\ar@{}[dr]|\square
    &
    N.
  }}
\end{equation}

\begin{equation}
    a_M\colon \cotb[M]\to\cotb[M],\quad(x;\xi)\mapsto(x;-\xi).
\end{equation}

同次シンプレクティック多様体\(\cotb[M]\)を考える．
すなわち，\(\cotb[M]\)の局所同次シンプレクティック座標\((x;\xi)\)に
おけるリウヴィル\(1\)形式は\[
    \alpha_M=\langle{\xi},dx\rangle
\]で与えられる．







\subsection{超局所台{\cite[5.1,6.5節]{KS90}}}
これもいつもの
\begin{DFN}\label{def11}
    \(\muS\)の定義
\end{DFN}
\begin{EG}\label{eg12}
    \(\muS\)の例
\end{EG}
\begin{THM}\label{thm13}
    包合性定理
\end{THM}
\subsection{局所化{\cite[6.1節]{KS90}}}
ここから知らない．

\(A\)を\(\cotb[M]\)の部分集合とし，
\(Z=\cotb[M]\smallsetminus A\)とする．
\(\muS(F)\subset Z\)をみたす層\(F\)からなる\(
    \Dompb(\kk_{M})
\)の充満部分圏\(\mucat[Z]{\kk_M}\)は三角圏である．
\[
    \Dompb(\kk_M;A)\coloneqq \Dompb(\kk_M)/\mucat[Z]{\kk_M}
\]
を\(\Dompb(\kk_M)\)の\(\mucat[Z]{\kk_M}\)による局所化とする．
すなわち，\(\Dompb(\kk_M;A)\)の対象は\(\Dompb(\kk_M)\)と同じだが，
\(\Dompb(\kk_M)\)の射\(u\colon F_1\to F_2\)については，
特三角\(F_1\to F_2\to F_3\overset{+1}{\longrightarrow}\)に
うめこんで\(\muS(F_3)\cap A=\varnothing\)となれば，
\(\Dompb(\kk_M;A)\)において同型になる．

\(p\in \cotb[M]\)で\(A=\{p\}\)とかけるとき，
\(\Dompb(\kk_M;\{p\})\)をたんに\(\Dompb(\kk_M;p)\)とかく．

\subsection{関手の演算{\cite[5.4節]{KS90}}}

\(M\)と\(N\)を実多様体とする．
\(q_i\)を\(M\times N\)からの射影，
\(p_i\)を\(\cotb[(M\times N)]\cong\cotb[M]\times \cotb[N]\)から
の射影とする．

\begin{DFN}\label{def14}
    非特性の定義
\end{DFN}

\(f\)が非特性的であることの\(\mdiff[f]\)による特徴づけ

\begin{THM}[{\cite[5.4節]{KS90}}参照]\label{thm15}
    \(f\colon M\to N\)を多様体の射とし，
    \(F\in\Dompb(\kk_M)\)，\(G\in\Dompb(\kk_N)\)とする．
    \begin{enumerate}[(i)]
        \item \(f\)が\(\supp(F)\)上固有とする．
        このとき，\(
            \muS(\Rder{f_{!}})
            \subset
            \mpb[f]\mdiff[f]^{-1}\muS(F)
        \)である.\label{item15-1}
        \item \(f\)が\(\muS(G)\)上非特性的であるとする．
        このとき，\(
            \muS(f^{-1}G)\subset\mdiff[f]\mpb[f]\muS(G)
        \)である.\label{item15-2}
        \item \(f\)が平滑であるとする．
        このとき，\(
            \muS(F)\subset M\mathbin{\times_{N}} \cotb[N]
        \)となることと，
        すべての\(j\in\zz\)に対し層\(H^j(F)\)が\(f\)の
        ファイバー上で局所定数層となることは同値である．\label{item15-3}
    \end{enumerate}
\end{THM}
\begin{CRL}\label{cor16}
    \(I\)を可縮な多様体とし，\(p\colon M\times I\to M\)を射影とする．
    \(F\in\Dompb(\kk_{M\times I})\)が\(
        \muS(F)\subset\cotb[M]\times\cotz[I]
    \)をみたすとき，\(F\cong p^{-1}\Rder{p_{\ast}}F\)となる．
\end{CRL}
\begin{proof}
    かく．
\end{proof}
\begin{CRL}\label{cor17}
    \(I\)を\(\rr\)の開区間とし，
    \(q\colon M\times I\to I\)を射影とする．
    \(F\in\Dompb(\kk_{M\times I})\)が\(
        \muS(F)\cap (\cotz[M]\times\cotb[I])
        \subset
        \cotz[(M\times I)]
    \)をみたし\(q\)が\(\supp(F)\)上固有であるとする．
    このとき，\(F_a\coloneqq \mres[F]{\{t=a\}}\)とおくと，
    任意の\(s,t\in I\)に対し，\(
        \RG(M;F_s)\cong \RG(M;F_t)
    \)が成り立つ．．
\end{CRL}
\begin{proof}
    次のカルテシアン図式を考える．
    \[
        \vcenter{\xymatrix@C=36pt@R=36pt{
        M\times\{s\}
        \ar[r]^-{a_M}
        \ar@{^{(}->}[d]^-{\iota_s}
        %\ar@{^{(}->}[r]_-{s_Y}&
        &
        \{s\}
        \ar@{^{(}->}[d]^-{\iota_s}\\
        M\times I
        \ar[r]^-{q}
        &
        I.
    }}\]
    \(q\)が固有なので\(\Rder{q_!}F\cong\Rder{q_{\ast}}F\)である．
    また，\(\mres[q]{M\times\{s\}}=a_M\)なので，\(a_M\)も固有である．
    したがって，
    \begin{align*}
        \RG(M;F_s)
        &\cong\Rder{{a_M}_{\ast}}\iota^{-1}F\\
        &\cong\Rder{{a_M}_{!}}\iota_{s}^{-1}F\\
        &\cong\iota_{s}^{-1}\Rder{{q}_{!}}F\\
        &\cong\iota_{s}^{-1}\Rder{{q}_{\ast}}F\\
        &\cong(\Rder{{q}_{\ast}}F)_{s}
    \end{align*}
    である．
    定理\ref{thm15}~\eqref{item15-2}から，\(
        \muS(\Rder{q_{\ast}}F)
        \subset 
        \mpb[q]\mdiff[q]^{-1}\muS(F)
    \)であり，
    仮定から
    \begin{align*}
        \mpb[q]\mdiff[q]^{-1}\muS(F)
        &=\mpb[q]((\cotz[M]\times\cotb[I])\cap\muS(F))\\
        &\subset \mpb[q](\cotz[(M\times I)])\\
        &= \cotz[I]
    \end{align*}
    である．
    よって，\(\muS(\Rder{q_{\ast}}F)\subset\cotz[I]\)となる．
    このとき\(V\in\Dompb(\kk)\)で\(\Rder{q_{\ast}}F\cong V_I\)となるものが存在する．
    \(t\)についても\(\Rder{q_{\ast}}F\cong V_I\)となるので，
    求める結果を得る．
\end{proof}





\subsection{モースの補題とモース不等式{\cite[5.4節]{KS90}}}

この節では\(C^1\)級関数\(\psi\colon M\to\rr\)を考える．
\begin{equation}\label{eq14}
    \Lambda_{\psi}=\left\{(x;d\psi(x))\right\}\subset\cotb[M]
\end{equation}
とおく．
以下の命題は超局所モースの補題（\cite[系5.4.19]{KS90}参照）の
特別な場合であり定理\ref{thm15}~\eqref{item15-2}から従う．
古典的な理論は定数層\(F=\kk_M\)に対応する．
\begin{PRP}\label{prp18}
    \(F\in\Dompb(\kk_M)\)とし，
    \(\psi\colon M\to\rr\)を\(C^1\)級関数で，
    \(\supp(F)\)上固有なものとする．
    \(a<b\)を実数とする．
    \(a\leqq \psi(x)<b\)に対し\(d\psi(x)\notin\muS(F)\)である
    とする．
    \(t\in\rr\)に対し，\(M_t=\psi^{-1}(]-\infty,t[)\)とおく．
    このとき，制限射\(\RG(M_b;F)\to\RG(M_a;F)\)は同型である．
\end{PRP}
\begin{proof}
    \(\psi\)が固有なので\(
        \Rder{\psi_{!}}F\cong\Rder{\psi_{\ast}}F
    \)である．いま，
    \begin{align*}
        \RG(M_a;F)&\cong\RG({\left]-\infty,a\right[}\mathrel{;}\Rder{\psi_{\ast}}F),\\
        \RG(M_b;F)&\cong\RG({\left]-\infty,b\right[}\mathrel{;}\Rder{\psi_{\ast}}F)        
    \end{align*}
    である．定理\ref{thm15}~\eqref{item15-2}より，
    \[
        \muS(\Rder{\psi_{\ast}}F)
        \subset
        \mpb[\psi]\mdiff[\psi]^{-1}\muS(F)
    \]
    なので，仮定の\(d\psi(x)\notin\muS(F)\)から，
    \(a\leqq\psi(x)<b\)に対し，
    \[
        \left(\RG_{[a,b[}\Rder{\psi_{\ast}}F\right)_{\psi(x)}=0
    \]
    となる．
\end{proof}
\begin{CRL}\label{cor19}
    \(F\in\Dompb(\kk_M)\)とし
    \(\psi\colon M\to\rr\)を\(C^1\)級関数とする．
    \(\Lambda_\psi\)を\eqref{eq14}で定めたものとする．
    次の条件が成り立つとする．
    \begin{enumerate}[(i)]
        \item \(\supp(F)\)はコンパクトである．
        \item \(\RG(M;F)\ne0\).
    \end{enumerate}
    このとき，\(\Lambda_{\psi}\cap\muS(F)\ne\varnothing\)である．
\end{CRL}

この節の終わりまでと\ref{ssec44}節では，\(\kk\)を体とする．
古典的なモース不等式は
層に拡張される（\cite{ST92}と\cite[命題5.4.20]{KS90}参照）．
そこでの結果を簡単に復習する．

各コホモロジーが有限次元である\(\kk\)ベクトル空間の有界複体\(E\)に対し，
\[
    b_j(E)=\dim H^j(E),\quad
    b_{l}^{\ast}(E)=(-1)^{l}\sum_{j\leqq l}^{}(-1)^{j}b_j(E)
\]
とおく．
\(C^1\)級関数\(\psi\colon M\to\rr\)を考え\(\Lambda_\psi\)を
上と同様に定める．
\(\psi\)がなめらかとは仮定していないことに注意．
\(F\in\Dompb(\kk_M)\)とする．
\begin{equation}\label{eq15}
    \text{
        集合\(\Lambda_{\psi}\cap\muS(F)\)は有限で，
        \(\{p_1,\dots,p_N\}のようにかける\)
    }
\end{equation}
とする．
\begin{equation}\label{eq16}
    x_i=\pi(p_i),\quad
    V_i\coloneqq\left(
        \RG_{\{\psi(x)\geqq\psi(x_i)\}}(F)
    \right)_{x_i}
\end{equation}
とおき，次を仮定する．
\begin{equation}\label{eq17}
    \text{すべての\(i\in\{1,\dots,N\}\)と\(j\in\zz\)に対し，
    空間\(H^{j}(V_i)\)は有限次元である．}
\end{equation}
\[
    b_j(F)=b_j(\RG(M;F)),\quad
    b_{j}^{\ast}(F)=b_{j}^{\ast}(\RG(M;F))
\]
とおく．
\begin{THM}
    \(F\in\Dompb(\kk_M)\)とし，\(\supp(F)\)はコンパクトとする．
    さらに\eqref{eq15}と\eqref{eq17}を仮定する．
    このとき，任意の\(l\)に対し\begin{equation}
        b_{l}^{\ast}(F)\leqq
        \sum_{i=1}^{N}b_{l}^{\ast}(V_i)\label{eq18}
    \end{equation}
    が成り立つ．
\end{THM}
\eqref{eq18}から，任意の\(j\)に対し
\begin{equation}
    b_{l}^{\ast}(F)\leqq
    \sum_{i=1}^{N}b_{l}^{\ast}(V_i)\label{eq18}
\end{equation}
が成り立つことがただちに従う．
\subsection{核{\cite[3.6節]{KS90}}}

\(M_i\) (\(i=1,2,3\)) を多様体とする．
\(
    M_{ij}\coloneqq M_i\times M_j
\) (\(1\leqq i,j\leqq 3\)), \(
M_{123}=M_1\times M_2\times M_3\)と略記する．
\(q_i\)で射影\(M_{ij}\to M_i\)または射影\(M_{123}\to M_i\)を表し，
\(q_{ij}\)で射影\(M_{123}\to M_{ij}\)を表す．
同様に，
\(p_i\)で射影\(M_{ij}\to M_i\)または射影\(
    \cotb[M]_{123}\to \cotb[M]_i
\)を表し，
\(p_{ij}\)で射影\(\cotb[M]_{123}\to \cotb[M]_{ij}\)を表す．
また\(\cotb[M_2]\)上の対蹠写像と\(p_{12}\)を合成した
写像\(p_{12^a}\)も導入する．

\(\Lambda_1\subset \cotb[M_{12}]\), 
\(\Lambda_2\subset \cotb[M_{23}]\)とする．
\begin{equation}\label{eq110}
    \Lambda_1\circ \Lambda_2
    \coloneqq
    p_{13}(
        p_{12^a}^{-1}\Lambda_1\cap p_{23}^{-1}\Lambda_2
    )
\end{equation}
とおく．
核のたたみこみ演算を考える：
\begin{align*}
    \conv[M_2]\colon 
    \Dompb(\kk_{M_{12}})\times\Dompb(\kk_{M_{23}})
    &\to\Dompb(\kk_{M_{13}})\\
    (K_1,K_2)&\mapsto K_1\conv[M_2]K_2\coloneqq
    \Rder{{q_{13}}_{!}}(q_{12}^{-1}K_{1}\ltens[]q_{23}^{-1}K_{2}).
\end{align*}
\(\Lambda_i=\muS(K_i)\subset\cotb[M_{i,i+1}]\)とし，
次を仮定する．
\begin{equation}\label{eq111}
    \left\{
        \parbox{60ex}{
            (i) $q_{13}$は$q_{12}^{-1}\supp(K_1)\cap q_{23}^{-1}\supp(K_2)$上で固有であり,
            \\%[2ex]
            (ii) 
            $%\ba[t]{l}
            p_{12^a}^{-1}\Lambda_1\cap p_{23}^{-1}\Lambda_2
            \cap (T^{\ast}_{M_1}M_1\times T^{\ast}M_2\times T^{\ast}_{M_3}M_3)\\%[1ex]
            \hspace*{10ex} 
            \subset T^{\ast}_{M_1\times M_2\times M_3}(M_1\times M_2\times M_3).$
    }\right.        
\end{equation}
定理\ref{thm15}から，仮定\eqref{eq111}のもとで次が成り立つ．
\begin{equation}
    \muS(K_1\conv[M_2]K_2)\subset \Lambda_1\circ\Lambda_2.
\end{equation}
混同の恐れがないときは\(\conv[M_2]\)のかわりに\(\conv[]\)とかく．

核のたたみこみの相対版も用いる．
多様体\(I\)と\(K_1\in\Dompb(\kk_{M_{12}\times I})\), 
\(K_2\in\Dompb(\kk_{M_{23}\times I})\)に対し，
\begin{equation}
    K_1\mathbin{\mres[\circ ]{I}}K_2\coloneqq
    \Rder{{q_{13I}}_{!}}(q_{12I}^{-1}K_{1}\ltens[]q_{23I}^{-1}K_{2})
\end{equation}
とおく．ただし\(q_{ijI}\)は射影\(M_{123}\times I\to M_{ij}\times I\)である，
上述の結果は相対的な場合にも成り立つ．
すなわち，次の条件を仮定する．
\begin{equation}\label{eq114}
    \left\{
        \parbox{60ex}{
            (i) \(\supp(K_1)\times_{M_{2}\times I}\supp(K_2)\to M_{13}\times I\)は固有であり，
            \\%[2ex]
            (ii) 
            $p_{12^{a}I^{a}}^{-1}\Lambda_1\cap p_{23I}^{-1}\Lambda_2
            \cap (T^{\ast}_{M_1}M_1\times T^{\ast}M_2\times T^{\ast}_{M_3}M_3)\times\cotb[I]\\%[1ex]
            \hspace*{10ex} 
            \subset T^{\ast}_{M_1\times M_2\times M_3\times I}(M_1\times M_2\times M_3\times I).$
    }\right.        
\end{equation}
ただし\[
    p_{12^{a}I^{a}}\colon
    \cotb[M_1]\times\cotb[M_2]\times\cotb[M_3]\times \cotb[I]
    \to\cotb[M_1]\times\cotb[M_2]\times \cotb[I]
\]は\[
    p_{12^{a}I^{a}}(v_1,v_2,v_3,u)=(v_1,-v_2,-u)
\]で与えられる．
このとき次が成り立つ．
\begin{equation}
    \muS(K_1\mres[\circ]{I}K_2)
    \subset 
    \Lambda_1\mres[\circ]{I}\Lambda_2
    \coloneqq
    r_{13}(r_{12^a}^{-1}\Lambda_1\cap r_{23}^{-1}\Lambda_2).
\end{equation}
ただし，図式
\[\vcenter{\xymatrix@C=16pt@R=36pt{
    &
    {
        \cotb[M_1]\times \cotb[M_2]
        \times \cotb[M_3]\times 
        (\cotb[I]\pb[I]\cotb[I])
    }
    \ar[dl]_-{r_{12^a}}
    \ar[d]_-{r_{13}}
    \ar[dr]^-{r_{23}}\\
    {\cotb[M_1]\times \cotb[M_2]\times \cotb[I]}
    &
    {\cotb[M_1]\times \cotb[M_3]\times \cotb[I]}
    &
    {\cotb[M_2]\times \cotb[M_3]\times \cotb[I]}
    }}
\]
において，
\begin{itemize}
    \item \(r_{12^a}\)は\(
        p_{12^a}\colon 
        \cotb[M_1]\times \cotb[M_2]\times \cotb[M_3]
        \to
        \cotb[M_1]\times \cotb[M_2]
    \)と第1射影\(\cotb[I]\times_{I}\cotb[I]\to\cotb[I]\)で与えられ，
    \item \(r_{23}\)は\(
        p_{23}\colon 
        \cotb[M_1]\times \cotb[M_2]\times \cotb[M_3]
        \to
        \cotb[M_2]\times \cotb[M_3]
    \)と第2射影\(\cotb[I]\times_{I}\cotb[I]\to\cotb[I]\)で与えられ，
    \item \(r_{13}\)は\(
        p_{13}\colon 
        \cotb[M_1]\times \cotb[M_2]\times \cotb[M_3]
        \to
        \cotb[M_1]\times \cotb[M_3]
    \)と加法の写像\(\cotb[I]\times_{I}\cotb[I]\to\cotb[I]\)で与えられる．
\end{itemize}


\subsection{局所有界圏と層のはりあわせ}

\textbf{準界} (prestack) \(
    U\mapsto \Domp(\kk_U)
\)（\(U\)は\(M\)の開集合）は\textbf{界} (stack) ではないが，
以下の古典的な結果があるので用いる．
\begin{LMM}
    \(U_1\)と\(U_2\)を\(M\)の開集合とし，
    \(U_{12}\coloneqq U_1\cap U_2\)とおく．
    \(F_{i}\in\Domp(\kk_{U_i})\) (\(i=1,2\)) とし，
    同型\(
        \varphi_{12}\colon 
        \mres[F_{1}]{U_{12}}\cong \mres[F_{2}]{U_{12}}
    \)が存在するとする．
    このとき，\(F\in\Domp(\kk_{U_1\cup U_2})\)と同型\(
        \varphi_{i}\colon \mres[F]{U_{i}}\cong F_{i}
    \) (\(i=1,2\)) で\(
        \varphi_{12}=
        \mres[\varphi_{2}]{U_{12}}
        \circ
        \mres[\varphi_{1}]{U_{12}}^{-1}
    \)をみたすものが存在する．
    しかも，そのような組\(
        (F,\varphi_1,\varphi_2)
    \)は（一意には決まらない）同型の違いを除けば1つしかない．
\end{LMM}
\begin{proof}
    かく．
\end{proof}
\begin{DFN}
    \(F\in\Domp(\kk_M)\)が\textbf{局所有界} (locally bounded) であるとは，
    任意の相対コンパクト開集合\(U\subset M\)に対し\(\mres[F]{U}\in\Dompb(\kk_U)\)となることをいう．
    局所有界な対象からなる\(\Domp(\kk_M)\)の充満部分圏を\(\Domplb(\kk_M)\)で表す．
\end{DFN}

\(\Dompb(\kk_M)\)に対して定義された局所的な概念は，
\(\Domplb(\kk_M)\)の対象に拡張される．
とくに超局所台も\(\Domplb(\kk_M)\)の対象に対して定まる．
グロタンディークの演算のうち有界性を保つものは，
順像は場合によっては除いて，局所有界性も保つ．
ところが，
\(F\in\Domplb(\kk_M)\)と多様体の射\(
    f\colon M\to N
\)で\(\supp(F)\)上固有なものに対して，
\(
    \Rder{f_{\ast}}F\cong\Rder{f_{!}}F\in\Domplb(\kk_M)
\)が成り立つので，
この状況でも定理\ref{thm15}が成り立つ．



\subsection{量子化接触変換{\cite[7.2節]{KS90}}}

\subsection{関手\(\muhom\) {\cite[4.4,7.2節]{KS90}}}

\subsection{単純層{\cite[7.5節]{KS90}}}

\section{余法線の対角集合への変形}

\section{同次ハミルトンイソトピーの量子化}

\subsection{量子化の一意性と台}

\subsection{量子化の存在：``コンパクト''の場合}

\subsection{量子化の存在：一般の場合}

\subsection{超局所台の変形}

\section{分離不能性問題への応用}

\subsection{分離不能性：同次の場合}

\subsection{分離不能性：モース不等式}

\subsection{分離不能性：非負ハミルトンイソトピー}

\subsection{分離不能性：シンプレクティックの場合}\label{ssec44}

\appendix


\section{ハミルトンイソトピー}

\subsection{シンプレクティック同型の族}

\subparagraph*{完全の場合}

\subsection{錐形ラグランジュ部分多様体の族}

\subsection{変数の追加}



%===============================================
% 参考文献スペース
%===============================================
\begin{thebibliography}{20} 
  \bibitem[GKS12]{GKS12} St\'ephane Guillermou, 
  Masaki Kashiwara, Pierre Schapira, 
  \textit{Sheaf Quantization of Hamiltonian Isotopies and Applications to Nondisplaceability Problems}, 
  Duke. Math. J. 161, No. 2, pp.201--245, 2012.  
  \bibitem[GP74]{GP74} Victor Guillemin, Alan Pollack, 
    \textit{Differential Topology}, 
    Prentice-Hall, 1974.
  \bibitem[KS90]{KS90} Masaki Kashiwara, Pierre Schapira, 
    \textit{Sheaves on Manifolds}, 
    Grundlehren der Mathematischen Wissenschaften, 292, Springer, 1990.
  \bibitem[Hat89]{Hat89} 服部晶夫, 多様体 増補版, 岩波全書 288, 岩波書店, 1989.
  %\bibitem[Og02]{Og02} 小木曽啓示, 代数曲線論, 朝倉書店, 2022.
  \bibitem[Sch85]{Sch85} Pierre Schapira, 
    \textit{Microdifferential Systems in the Complex Domain}, 
    Grundlehren der Mathematischen Wissenschaften, 269, 
    Springer, 1985.
    \bibitem[ST92]{ST92} Pierre Schapira, Nobuyuki Tose, 
    \textit{Morse Inequalities for \(\rr\)-constructible Sheaves}, 
    Adv. Math. 93, pp.1--8, 1992. 
\end{thebibliography}

%===============================================




\end{document}
